\documentclass[11pt]{article}
\usepackage[T1]{fontenc}
\usepackage{amsmath,amssymb}
\usepackage{booktabs}
\usepackage{geometry}
\usepackage[hidelinks]{hyperref}
\geometry{margin=1in}

\title{Experiment 010: H1 Presentation-Stress Report}
\author{Generated by run.py}
\date{}

\begin{document}
\maketitle

\section{Audit Data}
This TeX file is a compact report generated from the same in-memory result object written to \href{results.json}{\texttt{results.json}}.  It is computational evidence only and does not prove presentation invariance.

\section{Conventions}
\begin{itemize}
\item Base field: \(\mathbb{Q}\).
\item Category: \texttt{nonunital Jordan algebras}.
\item Target algebra: \(J=k\{x,y\}/(x^2,xy,y^2)\).
\item Matrix convention: source rows, target coordinates.
\item Canonical relation basis: \((x^2,\ xy,\ y^2)\).
\item Minimal relation cells: \((r_{xx},\ r_{xy},\ r_{yy})\).
\end{itemize}

\section{Outcome Summary}
\begin{center}
\scriptsize
\begin{tabular}{p{0.42\linewidth}p{0.21\linewidth}rrc}
\toprule
Case & Outcome & \(\dim C_1\) & \(\dim H_1\) & Cert. \\
\midrule
\texttt{A\_minimal\_presentation} & \texttt{pass\_strong} & 3 & 3 & yes \\
\texttt{C\_redundant\_relation\_without\_syzygy} & \texttt{expected\_fail\_observed} & 4 & 4 & no \\
\texttt{D\_redundant\_relation\_with\_syzygy} & \texttt{pass\_strong} & 4 & 3 & yes \\
\texttt{F\_truncation\_depth\_diagnostic} & \texttt{pass\_strong} & 4 & 3 & yes \\
\texttt{B\_linear\_change\_g1} & \texttt{pass\_strong} & 3 & 3 & yes \\
\texttt{B\_linear\_change\_g2} & \texttt{pass\_strong} & 3 & 3 & yes \\
\texttt{B\_linear\_change\_g3} & \texttt{pass\_strong} & 3 & 3 & yes \\
\texttt{E\_multiple\_redundant\_relations\_m1} & \texttt{pass\_strong} & 4 & 3 & yes \\
\texttt{E\_multiple\_redundant\_relations\_m2} & \texttt{pass\_strong} & 5 & 3 & yes \\
\texttt{E\_multiple\_redundant\_relations\_m3} & \texttt{pass\_strong} & 6 & 3 & yes \\
\texttt{E\_multiple\_redundant\_relations\_m5} & \texttt{pass\_strong} & 8 & 3 & yes \\
\bottomrule
\end{tabular}
\end{center}

\section{MVP Diagnostics}
\begin{itemize}
\item Minimal presentation: outcome \texttt{pass\_strong}, \(\dim H_1=3\).
\item Naive redundant-relation control: outcome \texttt{expected\_fail\_observed}, \(\dim H_1=4\).  This is the planned false extra $H_1$ class.
\item Repaired redundant relation: outcome \texttt{pass\_strong}, \(\operatorname{rank}\partial_2^Q=1\), \(\dim H_1=3\).
\item Minimum passed: \texttt{true}.
\end{itemize}

\section{Extended Stress Tests}
\subsection{Linear Changes Of Generators}
\begin{center}
\begin{tabular}{llc}
\toprule
Case & Outcome & \(\operatorname{Sym}^2(g)\) check \\
\midrule
\texttt{B\_linear\_change\_g1} & \texttt{pass\_strong} & yes \\
\texttt{B\_linear\_change\_g2} & \texttt{pass\_strong} & yes \\
\texttt{B\_linear\_change\_g3} & \texttt{pass\_strong} & yes \\
\bottomrule
\end{tabular}
\end{center}

\subsection{Multiple Redundant Relations}
\begin{center}
\begin{tabular}{lrrc}
\toprule
Case & \(\operatorname{rank}\partial_2^Q\) & \(\dim H_1\) & Cert. \\
\midrule
\texttt{E\_multiple\_redundant\_relations\_m1} & 1 & 3 & yes \\
\texttt{E\_multiple\_redundant\_relations\_m2} & 2 & 3 & yes \\
\texttt{E\_multiple\_redundant\_relations\_m3} & 3 & 3 & yes \\
\texttt{E\_multiple\_redundant\_relations\_m5} & 5 & 3 & yes \\
\bottomrule
\end{tabular}
\end{center}

\section{Truncation Depth}
\begin{center}
\begin{tabular}{lrrr}
\toprule
Presentation model & \(N=1\) & \(N=2\) & \(N=3\) \\
\midrule
\texttt{minimal\_presentation} & 3 & 3 & 3 \\
\texttt{redundant\_without\_syzygy} & 4 & 4 & 4 \\
\texttt{redundant\_with\_syzygy} & 4 & 3 & 3 \\
\bottomrule
\end{tabular}
\end{center}

\section{Basis-Alignment Certificates}
For redundant-relation cases, the strong certificate checks \(\operatorname{im}(\partial_2^Q)=\ker(\rho)\) exactly over \(\mathbb{Q}\).  The naive control intentionally omits the syzygy, so the kernel of \(\rho\) is not killed and the false extra $H_1$ class remains.

\begin{center}
\small
\begin{tabular}{p{0.44\linewidth}rrc}
\toprule
Case & \(\dim C_2\) & \(\operatorname{rank}\partial_2^Q\) & \(\operatorname{im}\partial_2^Q=\ker\rho\) \\
\midrule
\texttt{C\_redundant\_relation\_without\_syzygy} & 0 & 0 & no \\
\texttt{D\_redundant\_relation\_with\_syzygy} & 1 & 1 & yes \\
\texttt{E\_multiple\_redundant\_relations\_m1} & 1 & 1 & yes \\
\texttt{E\_multiple\_redundant\_relations\_m2} & 2 & 2 & yes \\
\texttt{E\_multiple\_redundant\_relations\_m3} & 3 & 3 & yes \\
\texttt{E\_multiple\_redundant\_relations\_m5} & 5 & 5 & yes \\
\bottomrule
\end{tabular}
\end{center}

\section{Limitations}
This is a bounded low-degree diagnostic for one square-zero example. It supports H1 stability under the listed presentation stresses but does not prove presentation invariance.

\section{Interpretation}
Experiment 010 supports the stability of the computed H1 for the two-generator square-zero Jordan algebra under deterministic linear generator changes and controlled redundant-relation refinements. The expected-fail control reproduces the false extra H1 class that appears when a redundant relation is added without the corresponding syzygy. This supports, but does not prove, the identification of the stable H1 with Sym\^{}2(k\^{}2) in this bounded example.

\end{document}
